\section{Discussion}
\label{sec:discussion}

\paragraph{What the audit adds.}
The counterexample separates a structural certificate from the computation
that replaces the removed structure. Separate routes expose both sources in
one reconstruction; the other loses a source at the readout. The transfer
makes this difference explicit and complements boundary-aware
controls~\cite{park2026}. Its contribution is the audit protocol and the
evidence that the two levels can disagree, rather than a new graph
architecture or a universal information law.

\paragraph{What the criterion guarantees.}
Proposition~1 is a population statement for one noiseless subspace task.
It says whether the XOR label is decodable from transferred source states;
it does not promise finite-MLP accuracy or biological generality.

\paragraph{Information is not utilization.}
All nonzero $k=8$ doses preserve ideal label information, while finite
accuracy varies substantially. Those curves describe how the recorded
procedure handles different input mixtures within an information-preserving
regime. Rank-one transfer can retain the label, so full source recovery is a
stronger requirement than this classification task needs.

\paragraph{Validity boundary.}
The evidence comes from one male connectome, two overlapping $k{=}8$
extractions, a fixed two-step linear operator and synthetic task labels.
The CIFAR-10 assay tests source-specific readouts on four mixed codes, not
recognition accuracy. The exact population criterion assumes noiseless
continuous features; nonlinear propagation requires an input-dependent
analysis. The two dense $k{=}8$ bases match frozen source hashes, but
historical dose-instance hashes and bit-identical initializations are
unavailable. Reconstructed matrices therefore give a consistency result,
and the historical gap cannot be assigned to one editing rule. A separate
reverse-support image control finds no connectome-specific benefit.

\section{Conclusion}
\label{sec:conclusion}

Graph interventions should report the information that survives at the
readout alongside the structure they preserve. The task-information audit
uses the executed operator to expose source access, transfer geometry and
their consequences for an ideal task, then keeps finite-model accuracy as a
separate measurement. On matched MaleCNS graphs, this protocol reveals
different information under the same structural certificate and shows that
information-preserving edits can still produce different learned behavior.
The result is a sharper standard for intervention studies: state the
operator, inspect the information reaching the readout, and report which
level of the analysis supports each conclusion.
